\documentclass[11pt,a4paper]{article}
\usepackage[T1]{fontenc}
\usepackage{lmodern,amsmath,amssymb,amsthm,mathtools,booktabs,microtype,xcolor}
\usepackage[margin=26mm]{geometry}
\usepackage[colorlinks=true,linkcolor=blue!45!black,citecolor=blue!45!black,urlcolor=blue!45!black]{hyperref}
\newtheorem{theorem}{Theorem}[section]
\newtheorem{lemma}[theorem]{Lemma}
\newtheorem{proposition}[theorem]{Proposition}
\newtheorem{corollary}[theorem]{Corollary}
\theoremstyle{remark}\newtheorem{remark}[theorem]{Remark}
\newcommand{\Q}{\mathbb Q}\newcommand{\Z}{\mathbb Z}
\DeclareMathOperator{\Tr}{Tr}\DeclareMathOperator{\Nm}{N}
\setlength{\emergencystretch}{2em}
\title{Simultaneous valuation conditions and explicit quartics}
\author{}\date{9 November 2026}
\begin{document}\maketitle
\begin{abstract}
We give a valuation argument that replaces several multiplicative local
tests by one polynomial expression. In the elliptic construction over
$\Q(\sqrt2)$, powers already computed for the formal logarithm provide
the required terms without additional high-degree equations. Sixteen
local tests of three different orders are expressed by five or six
such polynomials. The complete coefficient tables specify quartics
in 301 and 351 total variables.
Under the stated arithmetic premises, sufficiently large members of the
301-variable family define $\Z$ in $\Q$, and an explicit member of the
351-variable family does so. The argument preserves the formal-depth
requirements, the scalar normalization, and the numerical sufficient
parameter for the explicit construction.
\end{abstract}

\section{Local tests with different valuation classes}
Let $v$ be a discrete valuation with value group $\Z$, and extend it
by $v(0)=+\infty$. The following elementary observation is useful when
several conditions must hold at the same place.

\begin{lemma}\label{lem:distinct}
Suppose $q_j\ne0$ and $v(q_j)$ is odd for $0\le j<s$. Put
\[
 T(\boldsymbol q;\boldsymbol D)
 =1+\sum_{j=0}^{s-1}q_j^{2^j}D_j^{2^{j+1}}.
\]
The finite valuations of the displayed summands are pairwise distinct.
Consequently $T\ne0$ and
\[
 v(T)=\min\bigl(0,\ 2^j(v(q_j)+2v(D_j)):0\le j<s\bigr)\le0.
\]
If $v(T)=0$, every nonconstant summand has positive valuation and
$T\equiv1$ in the residue field.
\end{lemma}
\begin{proof}
For a nonzero $D_j$, the integer $v(q_j)+2v(D_j)$ is odd.
Its product with $2^j$ is divisible by $2^j$ and not by $2^{j+1}$.
Two such integers with different indices cannot coincide, and neither
can be zero. The minimum valuation among the nonzero summands is
therefore attained exactly once. The ultrametric inequality is an
equality in this situation, proving every assertion. Zero values of
$D_j$ simply omit their summands.
\end{proof}

We use $K=\Q(\rho)$, $\rho^2=2$, and write $\sigma$ for conjugation.
At a prime $p>2$ the valuations on $K$ above $p$ are unramified and
are normalized by $v(p)=1$.

\begin{proposition}\label{prop:general}
Let $R,\delta\in\Q^\times$ and suppose
\[
 \nu=v_p(R)\text{ is odd},\qquad v_p(\delta)=0.
\]
For a positive integer $m$ define
\[
 q_m=R^{1-2m}\delta^{-2m-2}.
\]
Let $W\in K$ be a product of finitely many factors, each of the form
\[
 1+\sum_{j=0}^{s-1}q_{m_j}^{2^j}D_j^{2^{j+1}},\qquad D_j\in K,
\]
where each $m_j$ is a positive integer and the orders may vary
within a factor and between factors. Zero tests are allowed. Write
$W=w_0+w_1\rho$. Suppose rational numbers $T_0,e,b$ satisfy
\begin{equation}\label{eq:trace}
 w_0^2-2w_0T_0-2w_1^2=0,\qquad e\delta b=T_0^2,
 \qquad e\delta=R^{-1}.
\end{equation}
Assume also the pole obstruction: $b\ne0$ implies that $v_p(b)$
cannot be a negative odd integer. Then $\nu>0$, and there is one place
$v\mid p$ at which every test in every factor satisfies
\[
                            v(D_j)\ge m_j.
\]
In particular, no inverse witness for $w_0$ is needed in this assertion.
\end{proposition}
\begin{proof}
For each $m$, $v(q_m)=(1-2m)\nu$ is odd. Lemma~\ref{lem:distinct}
therefore gives nonzero factors with nonpositive valuation at every
place above $p$. Thus $W\ne0$. If $w_0=0$, the first equation in
\eqref{eq:trace} gives $w_1=0$, a contradiction. Since $K$ is a field,
$\Nm(W)\ne0$, and the same equation gives
\[
 T_0=\frac{\Nm(W)}{2w_0}
     =\frac1{\Tr(W^{-1})}\ne0.
\]
The second equation now yields
\[
                        b=\frac{R}{\Tr(W^{-1})^2}\ne0.
\]
These equalities have been obtained from the literal polynomial rows;
none assumes a deleted inverse condition.

Put $\ell=\max_{v\mid p}v(W)\le0$. The trace inequality gives
\[
 v_p(\Tr(W^{-1}))\ge-\ell,
 \qquad v_p(b)\le\nu+2\ell.
\]
The valuation of $b$ is odd, since it differs from $\nu$ by an even
integer. The pole obstruction therefore gives
\begin{equation}\label{eq:positive}
                         0<v_p(b)\le\nu+2\ell.
\end{equation}
Hence $\nu>0$ and $\ell>-\nu/2$.

Choose a place attaining $\ell$. The valuation of any factor of $W$
is at least $\ell$, because all factor valuations are nonpositive
and their sum is $\ell$. Lemma~\ref{lem:distinct} then shows that
every nonzero term, whose test is assigned order $m_j$, has valuation
greater than $-\nu/2$. Thus
\[
 2^j\bigl((1-2m_j)\nu+2v(D_j)\bigr)>-\nu/2,
\]
or
\[
 v(D_j)>\left(m_j-\frac12-\frac1{2^{j+2}}\right)\nu
       \ge m_j-\frac34.
\]
The last inequality uses $m_j\ge1$ and $\nu\ge1$. Integrality of the
valuation proves $v(D_j)\ge m_j$. The inequality holds trivially when
$D_j=0$.
\end{proof}

\section{The rational height detectors}
For the curve used in the construction put
\[
 d=3-\rho,\quad D=d^2=11-6\rho,\quad
 E:y^2=x^3-Dx,\quad f(Q)=1/x(Q),\quad z(Q)=-x(Q)/y(Q).
\]
At the identity both functions extend by zero in the formal
neighborhood. The point charts and rational logarithms are those
specified in \cite{Explicit370,Family317}. The logarithms are unchanged
here, including their numerator and denominator nonvanishing conditions.
Write $U_0=H$, $V_0=H-P$, and $Q_i=U_i-V_i$ for the fixed root
points. These are elliptic-curve group differences. The exact addition
rows and identity-inclusive point charts in the cited specifications
enforce these relations; no scalar multiple with a noninteger index
is used to define a point.

Each height detector uses a nonzero rational square $a$, an input
$t=a\lambda$, and
\[
 \delta=t-1,\qquad
 \Theta=\left\{0,-\frac12,-1,\frac1{2399},-\frac{32}{407}\right\},
 \qquad F(t)=\prod_{\theta\in\Theta}(t-\theta).
\]
Its contact radical $M(t)$ is the product of primes $p>76176$ for
which $v_p(t-\theta)$ is positive and odd for some $\theta\in\Theta$.
The weight equations give
\begin{equation}\label{eq:weight}
                         F(t)e=(a\delta)^5.
\end{equation}
On $\delta F(t)\ne0$ define
\[
 R=\frac{F(t)}{a^5\delta^6},\qquad
 e=\frac1{R\delta},\qquad
 q_m=\frac{e^{2m-1}}{\delta^3}.
\]
The last expression equals the $q_m$ of Proposition~\ref{prop:general}.

\begin{lemma}\label{lem:contact}
Every solution of \eqref{eq:weight} has $F(t)\ne0$. If $p\mid M(t)$,
then $\delta\ne0$, $v_p(\delta)=0$, and $v_p(R)$ is odd. If $\delta=0$,
then $M(t)=1$.
\end{lemma}
\begin{proof}
The five roots in $\Theta$ differ from $1$. Therefore $F(t)=0$
contradicts \eqref{eq:weight}, as $a\ne0$. At a prime greater than
$76176$, all nonzero differences between distinct elements of
$\Theta\cup\{1\}$ are units. A contact prime divides exactly one
factor of $F(t)$ to positive odd order and makes $\delta$ a unit.
Since $a$ is a square, $v_p(F(t))-5v_p(a)$ is odd. Finally all prime
factors in the reduced numerators and denominators of $1-\theta$
are at most $76176$, so $M(1)=1$.
\end{proof}

Define the sixteen tests as follows. Their required orders are part of
these definitions:
\[
 \begin{array}{c|c|c}
 \text{Notation}&\text{Test}&\text{Order}\\\hline
 A_i&f(U_i)\ (1\le i\le4)&1\\
 B_i&\sigma f(V_i)\ (0\le i\le4)&1\\
 C_1&L_1-L'_1&10\\
 I&L_1-\beta&10\\
 D_H&f(H)&40\\
 C_i&L_i-L'_i\ (2\le i\le4)&40\\
 G&L_3^2-L_2(L_2-2)L_3L_4+(L_2-2)^2L_4^2-1&40
 \end{array}
\]
Here $\beta=51X^2+2504$, and the $L_i,L'_i$ are the retained primary
and alternate ratios, with the latter already conjugated. For an
ordered list of tests $(D_0,\ldots,D_{s-1})$, denote by
\begin{equation}\label{eq:factor}
 \mathcal T(D_0,\ldots,D_{s-1})
   =1+\sum_{j=0}^{s-1}q_{m(D_j)}^{2^j}D_j^{2^{j+1}}
\end{equation}
the factor formed with their respective orders.

The two-detector system uses the following six lists:
\begin{equation}\label{eq:lists301}
 \begin{gathered}
 (D_H,C_1,B_0),\quad(C_2,I,B_2),\quad(C_3,A_2,B_3),\\
 (C_4,A_3,B_4),\quad(G,A_4),\quad(A_1,B_1).
 \end{gathered}
\end{equation}
The three-detector system uses these five lists:
\begin{equation}\label{eq:lists351}
 \begin{gathered}
 (D_H,C_1,B_1,B_0),\quad(C_2,I,B_2),\quad(C_3,A_2,B_3),\\
 (C_4,A_3,B_4),\quad(G,A_4,A_1).
 \end{gathered}
\end{equation}
Each of the sixteen tests occurs exactly once in either collection.
Every order-forty test occupies position zero. The detector uses the
product $W$ of these factors, followed by \eqref{eq:trace} and the
unchanged normalized pole equations. Proposition~\ref{prop:general}
supplies every required local test order at one common place over
any contact prime. When $M(t)=1$, the contact-radical divisibility
statements are vacuous; no rational formula for $R$ or $W^{-1}$ is
used. This includes $t=1$.

\section{Preservation of simultaneous completeness}
The arithmetic hypotheses are those in the specified source construction
\cite[Section~2]{Source}: the elliptic group and index-height assertions,
the contact-height inequality, the pole obstruction and pole-realization
theorem, and the simultaneous prime-pattern theorem with its quantitative
weighted selection hypothesis. The latter allows an arbitrary fixed odd
avoidance modulus and the prescribed conjugate reduction orders, fixed
finite-place conditions, and real approximation conditions simultaneously.
The weight lower bound and slowly decreasing separation parameter are
retained exactly as stated there. This section shows that the present
replacement requires no stronger assertion.

\begin{proposition}\label{prop:complete}
For each integer input and each parameter admitted by the scalar
construction, the simultaneous point and prime selection in
\cite{Explicit370,Family317} supplies witnesses for the modified
height detectors. The local depths and the numerical sufficient scalar
parameter are unchanged.
\end{proposition}
\begin{proof}
First fix the scalar parameter, integer labels, root points, and positive
rational square scales by the retained scalar and Lucas construction.
Its rational functions are defined on nonempty real and finite-place
neighborhoods of the prepared points. Fix compact smaller neighborhoods
avoiding their finitely many poles. All test functions occurring in
\eqref{eq:lists301}--\eqref{eq:lists351} are bounded on these real neighborhoods.
The additional functions are only powers of the same tests; they
introduce no new poles or exceptional point exclusions. The continuity
modulus and finite avoidance set can be fixed at this stage, before the
contact primes are chosen.

The source selection has an exact product description that controls
all finite primes, not only the selected contacts. Choose a fixed
positive integer $D_0$ clearing every $D_0^2\theta/a$, and write
$v_*=D_0^2w^2$, $\lambda=u/v_*$. The simultaneous pattern makes each
of the five distinct forms $u-(\theta/a)v_*$ a positive selected
prime. Consequently
\begin{equation}\label{eq:prime-product}
 R=\frac{\displaystyle\prod_{\theta\in\Theta}
                 \bigl(u-(\theta/a)v_*\bigr)}{v_*^5\delta^6}.
\end{equation}
The denominator is a rational square. The selected primes avoid its
support and the fixed scale and contact-difference denominators, by
the avoidance conditions. Thus $v_p(R)=1$ at the five selected
contacts of this detector and $v_p(R)$ is even at every other prime.
This is the displayed product formula in the retained comparison
construction's completeness proof in \cite{Glazer}; it also follows
directly by substituting $t=a u/v_*$ in the definition of $R$.

At a selected contact prime the construction has $v_p(R)=1$ and a
designated primary place at which each original order-$m$ test satisfies
$v(D_j)\ge m$. Thus $q_mD_j^2$ has positive valuation there.
Every nonconstant term of a modified factor is a positive
integer power of one of the expressions $q_{m_j}D_j^2$.
Consequently each modified factor is $1$ modulo the primary prime.
Thus $W\equiv1$ there.

At the opposite split place, Lemma~\ref{lem:distinct} shows that each
modified factor either has negative valuation or is $1$ modulo the
prime. If at least one factor has negative valuation, then $W^{-1}$ is zero
modulo that prime; otherwise $W^{-1}\equiv1$. The two contributions
to $\Tr(W^{-1})$ therefore sum to $1$ or $2$ modulo $p$.
Since $p>76176>2$, the trace is a unit. Thus the pole input
$b=R/\Tr(W^{-1})^2$ has the required positive odd valuation at
every selected contact prime. The selection excludes the fixed
denominators of the root data and the scales, as in the source.

At the finitely many prepared places, the source preparation makes
each $q_mD_j^2$ as small as needed. Taking its second, fourth, or eighth power
preserves and strengthens this property. Hence the modified $W$ is
arbitrarily close to $1$ there. The required local squareclasses and
valuation bounds for $b$ at $2,3,17$ persist after increasing the
same finite preparation depths. At every prime outside the selected
contacts, \eqref{eq:prime-product} gives an even valuation of $R$,
and hence an even valuation of $b=R\Tr(W^{-1})^{-2}$, regardless
of its sign. At the selected contacts the preceding unit argument
gives valuation one. These two assertions together prove that the
reduced denominator of $b$ is a square. Moreover each selected prime
is square in $\Q_2,\Q_3,\Q_{17}$ by its prescribed congruences, so
\eqref{eq:prime-product} gives the required local squareclasses.

For real positivity and size, let $\lambda$ tend to positive infinity
after the root data and scales have been fixed. Uniformly on the chosen
compact real neighborhoods,
\[
 R=\frac{1+O(\lambda^{-1})}{a^6\lambda}>0,
 \qquad
 q_m=a^{10m-8}\lambda^{-3}(1+O(\lambda^{-1})).
\]
All new factors are positive at both real embeddings, because every $q_m>0$
and all powers of the tests are even. They tend uniformly to $1$;
a term in position $j$ decays as $O(\lambda^{-3\cdot2^j})$. Therefore
\[
 \Tr(W^{-1})\longrightarrow2,\qquad
 b=\frac{R}{\Tr(W^{-1})^2}\longrightarrow0^+.
\]
For sufficiently large $\lambda$, the pole input lies in $(0,17)$.
The same weighted selection premise supplies these arbitrarily large
parameters after the fixed congruences and avoidance conditions;
its decay requirement has not increased. Applying the pole-realization
input gives fresh simultaneous pole witnesses.

All remaining circuit coordinates are filled by their rational
definitions. In particular no assertion is made that an arbitrary
old pole witness is preserved. The numerical lower bound on the fixed
scalar parameter is used solely in soundness, where the test orders
and height inequalities are unchanged. Thus it is not enlarged by
this completeness argument.
\end{proof}

\section{A quadratic implementation using existing powers}
The variable $w=Df^2$ is used in the rational logarithm at each
$V_i$. For $i=0,2,3,4$, that part of the circuit already contains
$w^2$ and $w^4$. At $V_1$ it already contains $w^2$. These coordinates
are still required by the unchanged logarithm approximation.

The stored normalization is important. If $u,v$ are the two stored
coordinates for $w^2$, then the actual value is
\[
                           w^2=193u-132v\rho.
\]
Since $D^2=193-132\rho$ and $\Nm(D^2)=2401$, one has
\begin{equation}\label{eq:fourth}
 f^4=\frac{(37249u-34848v)+(25476u-25476v)\rho}{2401}.
\end{equation}
If $r,s$ are the stored coordinates for $w^4$, its actual value is
$r+s\rho$. As $D^4=72097-50952\rho$ and
$\Nm(D^4)=5764801$,
\begin{equation}\label{eq:eighth}
 f^8=\frac{(72097r+101904s)+(50952r+72097s)\rho}{5764801}.
\end{equation}
Conjugating the right sides gives the expressions consumed in
\eqref{eq:lists301}--\eqref{eq:lists351}. Thus all powers of the $V_i$ needed
there are linear forms in existing coordinates.

For \eqref{eq:lists301}, exactly five additional field powers are
needed: $C_1^4,I^4,A_2^4,A_3^4,A_4^4$. Each is the square of its
existing square. All four eighth powers $B_0^8,B_2^8,B_3^8,B_4^8$
are already available from \eqref{eq:eighth}; $B_1^4$ is available
from \eqref{eq:fourth}.

For \eqref{eq:lists351}, exactly nine additional field powers are
needed:
\[
 C_1^4,I^4,A_2^4,A_3^4,A_4^4,\quad
 B_1^8,B_0^{16},\quad A_1^4,A_1^8.
\]
The sixteenth power $B_0^{16}$ is the square of the existing eighth
power, and $B_1^8$ is the square of the existing fourth power.
Every listed field square uses two rational coordinates and two
quadratic equations. Explicitly, if $u+v\rho$ is a linear form in
available coordinates, its square is represented by
\[
                  j_0=u^2+2v^2,\qquad j_1=2uv.
\]
These equations have unique rational solutions for their output
coordinates, including when the field value is zero.

The affine normalizations of the input and Lucas squares are retained.
In source coordinates, they are
\[
 \begin{aligned}
 I^2&=(x_{307}-6x_{232}+306x_{248}+15015)
            +(2x_{308}-6x_{233})\rho,\\
 G^2&=(x_{319}-2x_{253}+1)+(4x_{320}-4x_{254})\rho.
 \end{aligned}
\]
The independent checker derives the same expressions and verifies
that the retained square rows have these interpretations.

The weight circuit stores $q=q_1$, $q_2=q^2$, and $q_4=q_2^2$.
Its retained quadratic identities on the nonzero domain are
\[
 a_1=q/e,\quad b_1=a_1^2/e,\quad c_1=b_1^2/q,\quad
 c_1d_1=qq_4,\quad c_1q_{10}=d_1^2,\quad c_1q_{40}=q_{10}^4.
\]
Since $q=e/\delta^3$, these yield
\[
 c_1=e^{-3}\delta^{-9},\quad d_1=e^8\delta^{-6},\quad
 q_{10}=e^{19}\delta^{-3},\quad q_{40}=e^{79}\delta^{-3}.
\]
The square of $q_{10}$ is already stored for the last identity. It
also supplies the order-ten terms in position one. Every order-forty
test occupies position zero, so no square or cube of $q_{40}$ is
required. For \eqref{eq:lists351}, add one scalar coordinate per
detector for $q_8=q_4^2$; this supplies the position-three term.
For \eqref{eq:lists301} no such coordinate is required. All weight
rows are quadratic. At a consumed contact prime their nonzero
interpretation follows from Lemma~\ref{lem:contact}; no division
is asserted at an empty contact radical.

These choices give the following exact coordinate accounting relative
to the specified systems with the same height scales:
\begin{center}
\begin{tabular}{lrr}
\toprule
Contribution&Two detectors&Three detectors\\\midrule
Starting coordinates&317&370\\
Shared expanded coefficients removed&$-22$&$-22$\\
New field powers&$+10$&$+18$\\
Squares and cubes of $q_{40}$ removed&$-4$&$-6$\\
Eighth powers of $q_1$&$0$&$+3$\\
Factor and product coordinates removed&$0$&$-12$\\\midrule
Total&301&351\\\bottomrule
\end{tabular}
\end{center}
There is a corresponding change in the number of quadratic rows.
The literal systems contain 275 and 323 rows respectively. All
remaining rows, including the cleared trace and the scalar guard,
are retained without alteration.

\section{Specified quartics and verification}
The exact tables specify primitive integer quadratic rows; their squared
sum is the defining polynomial. Every count includes the input coordinate.
The parameter $A=80989^{2N+1}$ is a fixed integer coefficient, not a
polynomial variable.
\begin{center}
\begin{tabular}{lrrr}
\toprule
Family&Total variables&Quadratic rows&Ordinary monomials\\
\midrule
Two height scales&301&275&4305\\
Three height scales&351&323&4990\\
\bottomrule
\end{tabular}
\end{center}
The ordinary monomial counts hold for $N\ge27464$. Before terms with
different powers of the fixed parameter are collected, the tables have
4320 and 5005 terms respectively. Every polynomial has degree exactly
four and content one. Both coefficient heights satisfy
\[
                             H(P_N)<2^{1520}80989^{4N+2}.
\]
The 351-variable member with $N=27464$ has a largest coefficient with
1791357 binary digits. The largest coefficient is its constant term.

The producer \nolinkurl{build_quartic_mixed.py} constructs both complete
systems and coefficient tables from the pinned source systems. The separate
program \nolinkurl{check_quartic_mixed.py} imports no producer. It
reconstructs the field constants from $D=11-6\rho$, verifies every changed
and unchanged quadratic row, and independently expands every square using
unordered pairs of terms. It checks the entire coefficient table, the
coordinate incidence and counts, the degree and content, and the exact
height of the explicit member. Its finite valuation checks supplement
Proposition~\ref{prop:general}; the proposition proves the assertion for
all odd valuations, without a search cutoff.

For completeness, the coefficient-height bound is a direct integer
certificate. For each coordinate monomial, write its coefficient as
$c_0+c_1A+c_2A^2$. The maximum of $|c_0|+|c_1|+|c_2|$ has 1520 binary
digits. Since $A\ge1$, the displayed bound follows. For $N\ge27464$,
$A>2^{1520}+1$, so the highest nonzero term of each such coefficient
dominates the sum of the lower terms. This also verifies that no ordinary
monomial disappears at an allowed parameter.

The conditional projection statements and their different effectiveness
properties are stated separately in \cite{Family301,Explicit351}.
Finite algebraic verification establishes none of the deep arithmetic
premises listed in Section~3.

\begin{thebibliography}{9}
\bibitem{Source}
\emph{A forced prime condition and an explicit quartic in 392 variables},
8 November 2026, Section~2; accompanying source table and independent
arithmetic certificates.
\href{../quartic/quartic392_20261108_v8.pdf}{PDF}; source: \nolinkurl{quartic392_20261108_v8.tex}.
\bibitem{Family317}
\emph{A family of quartics in 317 variables from aggregate trace heights},
8 November 2026, exact family and conditional projection theorem.
\href{../quartic_existence/quartic317_family_20261108_v9.pdf}{PDF}; source: \nolinkurl{quartic317_family_20261108_v9.tex}.
\bibitem{Explicit370}
\emph{A quartic in 370 variables with specified integer coefficients},
8 November 2026, exact family and simultaneous completeness proof.
\href{../quartic/quartic370_20261108_v10.pdf}{PDF}; source: \nolinkurl{quartic370_20261108_v10.tex}.
\bibitem{Family301}
\emph{A family of quartics in 301 variables}, 9 November 2026,
accompanying theorem and coefficient tables.
\href{../quartic_existence/quartic301_20261109_v11.pdf}{PDF}; source: \nolinkurl{quartic301_20261109_v11.tex}.
\bibitem{Explicit351}
\emph{An explicit quartic in 351 variables}, 9 November 2026,
accompanying theorem and coefficient tables.
\href{../quartic/quartic351_20261109_v11.pdf}{PDF}; source: \nolinkurl{quartic351_20261109_v11.tex}.
\bibitem{Glazer}
E.~Glazer, \emph{An explicit quartic defining the integers in the rationals},
Version~13, manuscript and coefficient data, commit
\nolinkurl{5bbd74f6a479776a332cbe036809d5dddae63720},
\url{https://github.com/elliotglazer/diophantine-z-q}.
\end{thebibliography}
\end{document}
